\BourbakiExerciseGroup

\begin{enumerate}
\def\labelenumi{\arabic{enumi})}
\tightlist
\item
  Let \(\mathbf{K}\) be a field of characteristic \(p > 0\) , E a separable extension of \(\mathbf{K}\) , of finite transcendence degree over \(\mathbf{K}\) .
\end{enumerate}

\begin{enumerate}
\def\labelenumi{\alph{enumi})}
\item
  If there exists a transcendence basis \(B_{0}\) of E over K and an integer \(m \geqslant 0\) such that \(\mathrm{K}(\mathrm{E}^{p^{m}})\) is separable over \(\mathrm{K}(\mathrm{B}_{0})\) , show that for every transcendence basis B of E over K there exists an integer \(n \geqslant 0\) such that \(\mathrm{K}(\mathrm{E}^{p^{n}})\) is separable over \(\mathrm{K}(\mathrm{B})\) .
\item
  Deduce that if the condition in a) holds, E admits a separating transcendence basis over K (if S is a p-basis of E over \(\mathrm{K}(\mathrm{E}^{p})\) and B a transcendence basis of E over K containing S (V, p.~176, Ex. 6), show that B is a separating transcendence basis of E over K, using Ex. 6 of V, p.~176).
\item
  Let K be a field of characteristic p \textgreater{} 0 and x an element transcendental over K (in an algebraically closed extension \(\Omega\) of K). Show that the union E of the extensions \(\mathrm{K}(x^{p^{-n}})\) of K is a separable extension of K such that tr. \(\deg_{K}E = 1\) but which does not admit a separating transcendence basis.
\item
  Let K be a perfect field. Show that the perfect closure of the field of rational fractions \(\mathrm{K}(\mathrm{T})\) is an extension of K of transcendence degree 1 which does not admit a separating transcendence basis.
\end{enumerate}

\begin{enumerate}
\def\labelenumi{\arabic{enumi})}
\setcounter{enumi}{1}
\item
  Let K be field of characteristic p \textgreater{} 0 and E an extension of finite transcendence degree over K. For E to admit a separating transcendence basis over K it is necessary and sufficient that E should be separable over K and the degree of imperfection of E over K equal to its transcendence degree over K.
\item
  Let \(\mathbf{E}\) be an extension of a field \(\mathbf{K}\) and \(\mathbf{F}\) an extension of \(\mathbf{E}\) .
\end{enumerate}

\begin{enumerate}
\def\labelenumi{\alph{enumi})}
\item
  If E admits a separating transcendence basis over K and if F admits a separating transcendence basis over E, then F admits a separating transcendence basis over K.
\item
  If F admits a separating transcendence basis over K and if tr. \(\deg_{K}E<+\infty\) , then E admits a separating transcendence basis over K (reduce to the case where tr. \(\deg_{K}F<+\infty\) and apply Ex. 1).
\item
  If F admits a finite separating transcendence basis over K and is separable over E then F admits a separating transcendence basis over E (use Ex. 1).
\end{enumerate}

\begin{enumerate}
\def\labelenumi{\arabic{enumi})}
\setcounter{enumi}{3}
\item
  Let K be a field of characteristic p \textgreater{} 0 whose absolute degree of imperfection (V, p.~170, Ex. 1) is equal to 1. Show that for an extension E of K to be separable over K it is necessary and sufficient that E should contain no element p-radical over K and not belonging to K (if \(a \in K\) forms a p-basis of K over \(K^{p}\) show that \(\mathbf{K}^{p^{-1}} = \mathbf{K}(a^{p^{-1}})\) and E are linearly disjoint over K).
\item
  \begin{enumerate}
  \def\labelenumii{\alph{enumii})}
  \tightlist
  \item
    Let K be a field of characteristic p \textgreater{} 0 and E an extension of K admitting a separating transcendence basis over K. Show that the intersection L of the fields \(\mathrm{K}(E^{p^{n}})\) , where n runs over N, is an algebraic extension of K. (Consider first the case where \(\operatorname{tr} \cdot \deg_{K} E < +\infty\) ; show that L is relatively perfect over K, hence E is separable over L. Now use Ex. 2 and Ex. 6 of V, p.~176. In the general case, if B is a separating transcendence basis of E over K, for each \(x \in L\) , there is a finite subset \(B_{0}\) of B such that the minimal polynomial of x over each of \(\mathrm{K}(B_{0}^{p^{n}})\) is the same as its minimal polynomial over \(\mathrm{K(B)}\) ; now deduce that if \(\mathrm{F} = \mathrm{K}(B_{0})(x)\) , then \(x \in \mathrm{K}(F^{p^{n}})\) for all n.)
  \end{enumerate}
\end{enumerate}

\begin{enumerate}
\def\labelenumi{\alph{enumi})}
\setcounter{enumi}{1}
\tightlist
\item
  Show that the largest perfect field \(E^{p^{\infty}}\) (intersection of the \(E^{p^{n}}\) ) contained in E is an algebraic extension of \(K^{p^{\infty}}\) (use a) and Ex. 8 of V, p.~176). Give an example where E is a finitely generated extension of a perfect field K but \(E^{p^{\infty}}\) is not contained in K.
\end{enumerate}

\begin{enumerate}
\def\labelenumi{\arabic{enumi})}
\setcounter{enumi}{5}
\tightlist
\item
  Let \(\mathbf{K}\) be a field of characteristic \(p > 0\) , \(\mathbf{E}\) an extension of \(\mathbf{K}\) such that there exists an integer \(h \geqslant 0\) for which \(\mathbf{K}(\mathrm{E}^{p^h})\) is separable over \(\mathbf{K}\) (which is always the case when \(\mathbf{E}\) is a finitely generated extension of \(\mathbf{K}\) ).
\end{enumerate}

\begin{enumerate}
\def\labelenumi{\alph{enumi})}
\item
  If C is a \(p\) -basis of E over \(\mathbf{K}(\mathbf{E}^p)\) , there is a subset B of C such that \(\mathbf{B}^{p^h}\) is a \(p\) -basis of \(\mathbf{K}(\mathbf{E}^{p^h})\) over \(\mathbf{K}(\mathbf{E}^{p^{h+1}})\) ; B is an algebraically free set over K. Show that the extension \(\mathbf{F} = \mathbf{K}(\mathbf{E}^{p^h})(\mathbf{B})\) is separable over K and E is an algebraic extension of F.
\item
  Show that B is a \(p\) -basis of F over \(\mathbf{K}(\mathbf{F}^p)\) . Show that if \(x \in \mathbf{E}\) and if \(m\) is the least integer such that \(x^{p^m} \in \mathbf{F}\) , then \(x^{p^m} \in \mathbf{K}(\mathbf{F}^{p^m})\) ; therefore E is a subfield of the field \(\mathbf{K}^{p^{-\infty}}(\mathbf{F})\) (isomorphic to \(\mathbf{K}^{p^{-\infty}} \otimes_{\mathbf{K}} \mathbf{K}\) ). A subextension F of E is said to be distinguished if it is separable over K and if E is a subfield of \(\mathbf{K}^{p^{-\infty}}(\mathbf{F})\) ; F is then a maximal separable extension of K in E.
\item
  Let P be a perfect field of characteristic p \textgreater{} 0, \(K = P(X, Y)\) the field of rational fractions in two indeterminates and \(\Omega\) an algebraically closed extension of the field of rational fractions \(\mathbf{K}(\mathbf{Z})\) over K. If \(u \in \Omega\) is such that \(u^{p} = X + YZ^{p}\) , show that \(\mathrm{E} = \mathrm{K}(\mathrm{Z}, u)\) is an extension of K such that \(\mathrm{K}(\mathrm{Z})\) and \(\mathrm{K}(u)\) are two distinguished subextensions of E, so that there exists no larger separable subextension of E. If \(v \in \Omega\) is such that \(v^{p^{2}} = X + YZ^{p}\) , show that in the extension \(\mathrm{E}' = \mathrm{K}(\mathrm{Z}, v)\) , \(\mathrm{K}(\mathrm{Z})\) is a maximal separable subextension of E but is not a distinguished subextension.
\end{enumerate}

\begin{enumerate}
\def\labelenumi{\arabic{enumi})}
\setcounter{enumi}{6}
\tightlist
\item
  Assume that the hypotheses of Ex. 6 hold, and further that the degree of imperfection of E over K (V, p.~170, Ex. 1) is finite (which is always the case if E is a finitely generated extension of K).
\end{enumerate}

\begin{enumerate}
\def\labelenumi{\alph{enumi})}
\tightlist
\item
  Let B be a \(p\) -basis of F over \(\mathbf{K}(\mathbf{F}^p)\) which is also a separating transcendence basis of F over K (and a transcendence basis of E over K). Show that for every integer \(k > 0\) , \(\mathbf{B}^{p^k}\) is a \(p\) -independent subset of \(\mathbf{K}(\mathbf{E}^{p^k})\) over \(\mathbf{K}(\mathbf{E}^{p^{k+1}})\) . Let \(q = \operatorname{tr} \cdot \deg_{\mathbf{K}} \mathbf{E}\) ; if \(m_k + q\) is the degree of imperfection of \(\mathbf{K}(\mathbf{E}^{p^k})\) over \(\mathbf{K}(\mathbf{E}^{p^{k+1}})\) , show that \(p^f\) , where \(f = \sum_{k=0}^{h-1} m_k\) , is equal to the degree {[}E : F{]} and is therefore independent of the choice of the
\end{enumerate}

distinguished subextension F of E. The integer \(p^{f}\) is called the order of inseparability of E over K and is written \([E:K]_{i}\) ; this number coincides with the inseparable degree denoted by \([E:K]_{i}\) in V, p.46 when E is of finite degree.

\begin{enumerate}
\def\labelenumi{\alph{enumi})}
\setcounter{enumi}{1}
\item
  For \(0 \leqslant k \leqslant h - 1\) , \([\mathbf{K}^{p^{-k}}(\mathbf{E}) : \mathbf{K}^{p^{-k}}]_i\) is equal to \(p^{f_k}\) , where \(f_k = \sum_{j=k}^{h-1} m_j\) .
\item
  Let \(L\) be a subextension of \(E\) , separable over \(K\) and such that \(E\) is \(p\) -radical and of finite degree over \(L\) ; show that \([E:L] \geqslant [E:K]_i\) (if \([L(E^{p^k}):K(E^{p^k})] = p^{r_k}\) , show that \(r_{k+1} \leqslant r_k + q\) , by considering a basis of \(L(E^{p^k})\) over \(K(E^{p^k})\) and a basis of \(L\) over \(K(L^p)\) ).
\item
  Let \(K'\) be an extension of K such that \(K'\) and E are contained in the same algebraically closed extension \(\Omega\) of K. Show that if \(\mathrm{K}'(\mathrm{F})\) is separable over \(K'\) , it is a distinguished subextension of the extension \(\mathrm{K}'(\mathrm{E})\) of \(K'\) . This is so when \(K'\) and E are algebraically disjoint over K and then we have \([\mathrm{K}'(\mathrm{E}):\mathrm{K}']_{i}\leqslant[\mathrm{E}:\mathrm{K}]_{i}\) ; equality holds when \(K'\) and E are linearly disjoint over K or when \(K'\) is a separable algebraic extension of K. e) Show that \([\mathrm{E}:\mathrm{K}]_{i}\) is the least value of \([\mathrm{E}:\mathrm{K}(\mathrm{B})]_{i}\) for all transcendence bases B of E over K; if \(\bar{K}\) is the relative algebraic closure of K in \(\Omega\) (which is an algebraic closure of K), then \([\mathrm{E}:\mathrm{K}(\mathrm{B})]_{i}=[\mathrm{E}:\mathrm{K}]_{i}[\bar{\mathrm{K}}(\mathrm{E}):\bar{\mathrm{K}}(\mathrm{B})]_{i}\) . If L is a subextension of E such that E is an algebraic extension of L and \(\bar{\mathrm{K}}(\mathrm{E})\) a separable (algebraic) extension of \(\bar{\mathrm{K}}(\mathrm{L})\) , then the largest separable extension F of L contained in E is distinguished.
\item
  Suppose that E is a finitely generated extension of K; show that for every subextension L of E we have \([L:K]_{i} \leqslant [E:K]_{i} \leqslant [L:K]_{i}[E:L]_{i}\) . Give an example where
\end{enumerate}

\[
[ \mathrm{L}: \mathrm{K} ] _ {i} = [ \mathrm{E}: \mathrm{K} ] _ {i} \quad \text { and } \quad [ \mathrm{E}: \mathrm{L} ] _ {i} > 1.
\]

\begin{enumerate}
\def\labelenumi{\arabic{enumi})}
\setcounter{enumi}{7}
\tightlist
\item
  Let \(k\) be a field and \(L\) an extension of \(k\) . Show that the following conditions are equivalent:
\end{enumerate}

\begin{enumerate}
\def\labelenumi{(\roman{enumi})}
\item
  L is a separable algebraic extension of a finitely generated pure transcendental extension.
\item
  \(L\) is a separable extension of \(k\) and
\end{enumerate}

\[
[ \Omega_ {L / k}: L ] = \operatorname{tr} \cdot \deg_ {k} L <   + \infty .
\]

(To prove (ii) \(\Rightarrow\) (i) we may proceed as follows: let \(x_{1},\ldots ,x_{n}\) in L be such that \(dx_{1},\ldots ,dx_{n}\) form a basis of \(\Omega_k(\mathbf{L})\) and put \(\mathbf{K} = k(x_1,\dots,x_n)\) . One first shows that \(\operatorname {tr.deg}_k(\mathbf{K}) = n\) and hence that K is a pure transcendental extension of \(k\) . Let \(\mathbf{L}'\) be a finite algebraic extension of K contained in L. By comparing the canonical mapping \(u:\Omega_k(\mathbf{K})\otimes_\mathbf{K}\mathbf{L}'\to \Omega_k(\mathbf{L}')\) with the canonical mapping \(\Omega_k(\mathbf{K})\otimes_\mathbf{K}\mathbf{L}\to \Omega_k(\mathbf{L})\) , show that \(u\) is injective. Calculate \([\Omega_k(\mathbf{L}'):\mathbf{L}']\) (V, p.~134, Cor. 1). Deduce that \(u\) is bijective and hence that \(\Omega_{\mathbf{K}}(\mathbf{L}') = 0\) . Show that \(\mathbf{L}'\) is separable over K; deduce that L is separable over K.)

\begin{enumerate}
\def\labelenumi{\arabic{enumi})}
\setcounter{enumi}{8}
\tightlist
\item
  With the notation of Ex. 8 let \(\mathbf{M} \subset \mathbf{L}\) be a subextension. Put \(m = \operatorname{tr} \cdot \deg_k \mathbf{M}\) , \(n = \operatorname{tr} \cdot \deg_k \mathbf{L}\) , \(r = n - m\) . Let \(x_1, \ldots, x_n\) in \(\mathbf{L}\) be such that \(\mathbf{L}\) is a separable algebraic extension of \(k(x_1, \ldots, x_n) = \mathbf{K}\) .
\end{enumerate}

\begin{enumerate}
\def\labelenumi{\alph{enumi})}
\item
  Show that, after renumbering the \(x_{i}\) if necessary, we may suppose that \(L\) is algebraic over \(M(x_{1},\ldots ,x_{r})\) .
\item
  Show that there exists a finitely generated subextension \(\mathbf{N} \subset \mathbf{M}\) such that:
\end{enumerate}

\(\alpha)\) L is algebraic over \(\mathbf{N}(x_1,\dots,x_s)\) ;

\(\beta)\) \(\mathbf{K}' = \mathbf{K}(\mathbf{N}(x_1,\dots,x_r))\) is a finite extension of \(\mathbf{N}(x_1,\dots,x_r)\) ;

\(\gamma)\) the canonical homomorphism \(\mathbf{K}^{\prime}\otimes_{\mathbf{N}(x_1,\ldots ,x_r)}\mathbf{M}(x_1,\dots,x_r)\to \mathbf{K}(\mathbf{M}(x_1,\dots,x_r)) = \mathbf{K}''\) is bijective.

\begin{enumerate}
\def\labelenumi{\alph{enumi})}
\setcounter{enumi}{2}
\tightlist
\item
  Let N be a subextension as in b). Show successively that \(K''\) is a separable algebraic extension of \(K'\) (because L is separable over K), that \(M(x_{1}, \ldots, x_{r})\) is a separable algebraic extension of \(N(x_{1}, \ldots, x_{r})\) (use b), \(\gamma)\) and that M is a separable algebraic extension of N. d) Show that M is a separable algebraic extension of a pure transcendental extension (use c) and the fact that N is a finitely generated separable extension of k).
\end{enumerate}

